\section{Анализ предметной области и формирование требований к гибридной интеллектуальной системе}

\subsection{Характеристика предметной области}
% Описываются границы предметной области, ключевые объекты, процессы,
% свойства, отношения, ограничения и участники.
%
% Следует зафиксировать терминологию, используемую далее в пояснительной
% записке, а также определить источники знаний и данных, используемые
% гибридной интеллектуальной системой.

\subsection{Анализ решаемого класса задач}
% Определяется класс задач, решаемых гибридной интеллектуальной системой:
% диагностика, классификация, выбор, конфигурирование, планирование,
% прогнозирование, объяснение, поиск знаний, формирование рекомендаций
% или иной класс задач.
%
% Описываются условия задачи, требуемые результаты, критерии корректности
% и качества решения, а также ситуации, в которых требуется объединение
% логико-символьного и нейробионического подходов.

\subsection{Анализ базы знаний, решателя задач и интерфейса}
% Описываются результаты предыдущих этапов курсового проектирования:
% - база знаний, разработанная в 4-м семестре;
% - решатель задач, разработанный в 5-м семестре;
% - интерфейс, разработанный в 6-м семестре.
%
% Для каждого компонента указываются назначение, используемые знания
% и данные, входы и выходы, ограничения, степень готовности и возможность
% включения в гибридную интеллектуальную систему.
%
% Определяется, какие компоненты используются без изменений, какие
% дорабатываются и какие реализуются заново.

\subsection{Анализ существующих гибридных интеллектуальных систем}
% Рассматриваются существующие гибридные интеллектуальные системы,
% нейросимвольные системы, RAG-системы, экспертные системы с компонентами
% машинного обучения, многоагентные системы и иные релевантные решения.
%
% Для каждого подхода указываются архитектура, способ интеграции компонентов,
% преимущества, ограничения и применимость к выбранной предметной области.

\subsection{Требования к гибридной интеллектуальной системе}
% Формулируются функциональные требования:
% - поддерживаемые классы задач;
% - поддерживаемые пользовательские и системные запросы;
% - использование базы знаний, решателя задач и интерфейса;
% - формирование решения, объяснения, рекомендаций или сообщений об ошибках;
% - обработка отсутствующих, неполных, противоречивых или некорректных данных.
%
% Формулируются нефункциональные требования:
% - корректность, полнота, устойчивость и объяснимость;
% - производительность и допустимое время отклика;
% - модульность, расширяемость, сопровождаемость;
% - возможность независимого обновления базы знаний, решателя задач,
%   интерфейса и компонента нейробионического подхода.
%
% Данный подраздел завершает аналитическую главу. В нём не следует
% описывать архитектурные решения, алгоритмы реализации, программный код
% или конкретные библиотеки.

\subsection{Выводы}
% Кратко фиксируются результаты анализа:
% характеристики предметной области и класса задач, состояние исходных
% компонентов, обоснование гибридного подхода и сформированные требования
% к гибридной интеллектуальной системе.
